\subsection*{Materiales}
\begin{itemize}
	\item Protoboard.
	\item Generador de voltaje DC con puesta a tierra.
	\item 11 resistencias \qty{1}{\kilo\ohm}$\pm 5\%$.
	\item Multímetro.
	\item Caimanes.
	\item Cables de conexión delgados.
\end{itemize}
\subsection*{Procedimiento}
En la práctica experimental se construyó un CDA utilizando la protoboard, resistencias
y cables sin polaridad.

Para realizar el montaje, utilizando la protoboard se colocaron las resistencias y los
cables de conexión como se indica en la Fig. \ref{fig:circuito}, tomando $2R$ como dos
resistencias en serie. Después, se estableción el voltaje del generador en
$\qty{16}{\volt}$ DC y se conectó en los lugares donde están las fuentes $V_0$ en la Fig.
\ref{fig:circuito}. Similarmente, se concectó la puesta a tierra del generador en el
lugar de $\Ground$ en el circuito. Por último, se conectó el multímetro entre los nodos
$d$ y $\Ground$.

Posteriormente, para medir las ocho posibles salidas del circuito $(0,...,7)\unit{\volt}$,
se encendió cada uno de los bits conectando las terminales con tres cables distintos; uno
para cada bit, y se registró el voltaje $V_\text{out}$ del voltímetro para cada combinación.